\documentclass[runningheads]{llncs}

\usepackage{graphicx}
\usepackage{amsmath,amssymb,amsfonts}
\usepackage{booktabs}
\usepackage{cite}
\usepackage{url}
\usepackage{hyperref}
\usepackage{xcolor}
\usepackage{algorithm}
\usepackage{algorithmic}

\begin{document}

\title{Safety-Aware Reinforcement Learning for Autonomous Mobile Robot Navigation: A Comparative Study of Predictive Safety Filtering and Recovery Mechanisms}
\titlerunning{Safety-Aware RL for Autonomous Mobile Robot Navigation}

\author{Author One\inst{1}\orcidID{0000-0000-0000-0000} \and
Author Two\inst{1} \and
Author Three\inst{1} \and
Faculty Advisor\inst{1}}

\authorrunning{A. One et al.}

\institute{Department of Computer Science \& Engineering / Robotics and Automation,\\
Institution Name, City, State, India\\
\email{\{author1, author2, author3, advisor\}@institution.edu.in}}

\maketitle

\begin{abstract}
Autonomous mobile robot (AMR) navigation in GPS-denied indoor environments has traditionally relied on classical receding-horizon planners (such as the Dynamic Window Approach, DWA) or model-free Deep Reinforcement Learning (DRL) local controllers. However, existing studies report only aggregate success rates, obscuring a critical underlying trade-off. In this paper, we uncover and mathematically formalize the \textbf{Safety vs. Conservatism Failure Mode Dichotomy}: under identical sensing and kinematic limits, classical DWA fails exclusively via conservative local-minima timeouts ($7.0\%$ on test, $14.0\%$ on ZigZag, $0.0\%$ collisions) due to obstacle inflation invalidations, whereas an unshielded Proximal Policy Optimization (PPO) controller fails exclusively via collisions ($16.0\%$ on test, $0.0\%$ timeouts) driven by saturated forward velocity ($v \to 0.26\,\text{m/s}$). On structurally unseen switchback labyrinths (held-out ZigZag family), unshielded PPO suffers a severe generalization collapse to $50.0\%$ success with a $50.0\%$ collision rate.

To eliminate collisions without inducing freezing deadlocks, we propose a novel \textbf{Medial-Axis Shielded Reinforcement Learning (MAS-RL)} architecture comprising: (1) a quadratic clearance-weighted geodesic field that mathematically forces global lookahead waypoints to track the medial axis of narrow passages; (2) an autonomous dynamic impasse detection engine that discovers alternative topological corridors via sub-$4\,\text{ms}$ wave-front replanning; and (3) an active 4-layer predictive safety shield providing an analytical discrete-time Control Barrier Function (CBF) stopping-distance governor ($v \le \sqrt{2 a_v^{\max} d_{\text{clear}}}$), flank-pinch protection, doorway centering torque, and non-sticky wall disengagement. Evaluated across 100 matched scenario trials with 95\% Wilson Score Confidence Intervals, MAS-RL achieves \textbf{90.0\% success [0.83, 0.94]} ($0.881$ SPL) on familiar layouts, and restores held-out ZigZag performance to \textbf{84.0\% [0.76, 0.90]} ($0.830$ SPL) with \textbf{0.0\% timeouts}, approaching DWA ($86.0\%$) while executing at $1.2\,\text{ms}$ on CPU ($7\times$ faster). We validate the framework in an interactive graphical testbed and demonstrate an embodied ROS 2 Jazzy, Gazebo Harmonic, and SLAM Toolbox deployment architecture operating at $10\,\text{Hz}$.

\keywords{Autonomous Mobile Robots \and Deep Reinforcement Learning \and Proximal Policy Optimization \and Dynamic Window Approach \and Control Barrier Functions \and SLAM Toolbox \and Navigation2 \and Medial Axis Transform.}
\end{abstract}

\section{Introduction}
Safe and efficient autonomous mobile robot (AMR) navigation in unstructured, GPS-denied indoor facilities (e.g., smart warehouses, hospitals, and search-and-rescue environments) requires concurrent spatial mapping, real-time localization, and collision-free local control~\cite{macenski2021slam, macenski2020marathon, macenski2023survey}. In modern robotics frameworks such as ROS 2 Navigation2 (Nav2)~\cite{macenski2020marathon}, navigation is hierarchically decoupled: a global planner computes a topological path over an occupancy grid, while a local controller generates high-frequency velocity commands $(v, \omega)$ to track that path while avoiding local obstacles detected by exteroceptive sensors.

Despite widespread commercial use, classical trajectory rollout planners such as the Dynamic Window Approach (DWA)~\cite{fox1997dwa} and Timed Elastic Bands (TEB) suffer from the notorious ``freezing robot'' problem: in tight doorways or cluttered passages, obstacle inflation layers frequently invalidate all forward candidate trajectories, causing the AMR to decelerate to zero velocity and trigger timeouts. Deep Reinforcement Learning (DRL)~\cite{schulman2017ppo, tai2017virtual} has emerged as an agile, data-driven alternative capable of sub-millisecond inference and fluid doorway traversal. However, existing DRL navigation literature~\cite{kastner2021arena, kastner2022arenabench, perumalla2020barn} suffers from three critical flaws:
\begin{enumerate}
    \item \textbf{Omission of Failure-Mode Analysis:} Studies focus on aggregate success rates, obscuring the physical divergence between classical conservatism and learned aggressiveness.
    \item \textbf{Out-of-Distribution (OOD) Generalization Cliff:} Policies trained on familiar rectilinear topologies experience catastrophic collision spikes when transferred to unseen labyrinthine switchbacks.
    \item \textbf{Absence of Statistical Rigor:} Results are frequently reported over small sample sizes ($N \le 10$) without confidence bounds, seed variance reporting, or closed-loop SLAM verification.
\end{enumerate}

To bridge these gaps, this paper establishes five core contributions:
\begin{enumerate}
    \item \textbf{Formalization of the Safety-Conservatism Dichotomy:} We prove experimentally across 100 matched trials that DWA achieves zero collisions ($0.0\%$) but fails strictly via conservative timeouts ($7.0\%$ test, $14.0\%$ ZigZag), whereas raw PPO achieves zero timeouts ($0.0\%$) but collapses on held-out labyrinths via collisions ($50.0\%$).
    \item \textbf{Medial-Axis Clearance-Weighted Planning \& Taut String-Pulling:} We formulate an 8-connected quadratic clearance penalty on the global geodesic field that naturally centers paths along corridor and doorway medians, coupled with $\Theta^*$-style string-pulling for true Euclidean shortcutting.
    \item \textbf{Autonomous Dynamic Impasse Replanning:} We formulate an impasse watchdog that dynamically detects stagnation, triggering sub-$4\,\text{ms}$ wave-front discovery of open alternative corridors.
    \item \textbf{Active 4-Layer Dynamic Safety Shield (Discrete CBF):} We derive an analytical stopping governor ($v \le \sqrt{2 a_v^{\max} d_{\text{clear}}}$), lateral repulsion, and non-sticky wall disengagement. This boosts nominal success to \textbf{90.0\% [0.83, 0.94]}, and restores held-out ZigZag success to \textbf{84.0\% [0.76, 0.90]} ($0.0\%$ timeouts).
    \item \textbf{Dual-Track Verification \& Interactive GUI:} Decoupling high-sample-efficiency 2D Gymnasium benchmarking ($>10{,}000$ steps/sec) from embodied ROS 2 Jazzy, Gazebo Harmonic, and SLAM Toolbox deployment ($1.2\,\text{ms}$ inference latency).
\end{enumerate}

\section{Related Work}
\textbf{Classical Planning and Nav2:} Classical local planners rely on dynamic velocity window sampling (DWA~\cite{fox1997dwa}) or non-linear graph optimization (TEB). Modern Nav2 stacks~\cite{macenski2020marathon, macenski2023survey} incorporate Model Predictive Path Integral (MPPI) control, sampling thousands of GPU rollouts. While capable, both DWA and MPPI remain vulnerable to inflation-layer deadlocks in constrictions.

\textbf{DRL Benchmarking Suites:} Francis et al.~\cite{francis2020prmrl} introduced PRM-RL, demonstrating that global topological graphs combined with RL local policies enable long-range indoor navigation. Wijmans et al.~\cite{wijmans2020ddppo} demonstrated decentralized distributed PPO for large-scale point-goal navigation. Kästner et al.~\cite{kastner2021arena, kastner2022arenabench} developed Arena-Rosnav and Arena-Bench to benchmark DRL in ROS/ROS 2, while Perumalla et al.~\cite{perumalla2020barn} introduced the BARN dataset for constrained spaces. Our work extends these benchmarks by reporting 95\% Wilson Score and bootstrap confidence intervals, formalizing failure modes, and demonstrating out-of-distribution recovery.

\textbf{Runtime Safety Shields and Barrier Functions:} Control Barrier Functions (CBFs)~\cite{ames2019cbf, cheng2019endtoend} ensure forward invariance of safe state sets via online Quadratic Programs (CBF-QPs). Execution-time action replacement~\cite{chiang2021execution} and Velocity Obstacle shields~\cite{kastner2022velocity} substitute recovery actions upon hazard detection. In this work, our analytical stopping-distance governor enforces discrete-time CBF guarantees with sub-millisecond execution overhead.

\section{System Formulation and Architecture}

\subsection{Kinematic State and Motion Model}
We evaluate a differential-drive ground robot conforming to the TurtleBot3 Waffle ($r_{\text{robot}} = 0.22\,\text{m}$, $v_{\max} = 0.26\,\text{m/s}$, $\omega_{\max} = 1.0\,\text{rad/s}$, $a_v^{\max} = 1.0\,\text{m/s}^2$, $a_\omega^{\max} = 3.2\,\text{rad/s}^2$). Over discrete interval $\Delta t = 0.1\,\text{s}$ (10 Hz), the unicycle state $\mathbf{p}_t = [x_t, y_t, \theta_t]^T$ transitions as:
\begin{align}
x_{t+1} &= x_t + v_t \cos(\theta_t) \Delta t, \quad y_{t+1} = y_t + v_t \sin(\theta_t) \Delta t \\
\theta_{t+1} &= \text{wrap}\left(\theta_t + \omega_t \Delta t\right)
\end{align}
subject to acceleration bounds $|v_{t+1} - v_t| \le a_v^{\max} \Delta t$ and $|\omega_{t+1} - \omega_t| \le a_\omega^{\max} \Delta t$.

\subsection{Markov Decision Process (MDP) Formulation}
The observation vector $\mathbf{s}_t \in \mathbb{R}^{29}$ comprises 24 angular sector range readings $\bar{r}_{0:23} \in [0, 1]$ (pooling a $360^\circ$ LiDAR scan over $15^\circ$ arcs, normalized by $R_{\max} = 3.5\,\text{m}$), normalized target distance $\bar{d}_g = \min(d_{\text{target}}/5.0, 1.0)$, target heading trigonometric features $[\sin(\phi_g), \cos(\phi_g)]$, and normalized velocities $[\bar{v}_t, \bar{\omega}_t]$. Continuous actions $\mathbf{a}_t = [a_v, a_\omega] \in [-1, 1]^2$ map to $v_{\text{cmd}} = \frac{a_v + 1.0}{2.0} v_{\max} \in [0, v_{\max}]$ (forward-only) and $\omega_{\text{cmd}} = a_\omega \omega_{\max}$.

The reward function guides the agent using a privileged 8-connected geodesic field $\Phi_{\text{geo}}$:
\begin{equation}
\mathcal{R}_t = 10.0 \cdot \left[ \Phi_{\text{geo}}(\mathbf{x}_{t-1}) - \Phi_{\text{geo}}(\mathbf{x}_t) \right] - 0.02 - 0.02 \frac{|\omega_t|}{\omega_{\max}} + R_{\text{prox}} + R_{\text{terminal}}
\end{equation}
where $R_{\text{prox}} = -0.5 (1.0 - \min(\mathbf{r}_t)/0.6)$ when $\min(\mathbf{r}_t) < 0.6\,\text{m}$, and $R_{\text{terminal}} \in \{+50.0 \text{ (goal)}, -50.0 \text{ (collision)}\}$.

\subsection{Medial-Axis Clearance-Weighted Planning \& Taut String-Pulling}
To prevent the global plan from clipping obstacle corners while preserving true Euclidean path optimality, Euclidean distance transform values $D_{\text{clear}}(r, c)$ relative to the inflated boundary are penalized quadratically below margin $d_{\text{margin}} = 0.25\,\text{m}$:
\begin{equation}
\text{Cost}_{\text{clear}}(r, c) = \left( \frac{d_{\text{margin}} - D_{\text{clear}}(r, c)}{d_{\text{margin}}} \right)^2 \cdot 0.35
\end{equation}
Edge traversal costs in Dijkstra wave-front propagation are scaled by $(1.0 + \text{Cost}_{\text{clear}})$. Setting $\alpha_{\text{clear}} = 0.35$ provides an essential Voronoi medial-axis tie-breaker: it naturally centers global paths along corridor and portal medians without inducing the metric distortions or multi-room detour artifacts caused by heavy clearance penalties. 

Furthermore, raw 8-connected grid paths are pruned via greedy line-of-sight shortcutting (\textit{Taut String-Pulling}, $\Theta^*$-style pruning). For discrete path coordinates $\{P_0, \dots, P_M\}$, intermediate nodes are pruned if collision-free clearance $\ge r_{\text{robot}} + 0.04\,\text{m}$ holds along the straight Euclidean chord, eliminating jagged staircases and yielding minimal-length Euclidean ribbons. When passing through doorways ($D_{\text{clear}} < r_{\text{robot}} + 0.12\,\text{m}$), lookahead is adaptively contracted ($L_{\text{eff}} = \min(L_{\text{nom}}, \max(0.65, 2.2(D_{\text{clear}} - r_{\text{robot}} + 0.04)))$) to enforce orthogonal portal entry before executing turning maneuvers into adjacent rooms.

\subsection{Safe Standoff Boundary Handling \& Physical Obstacle Replanning}
In practical warehouse deployment, navigational destinations are frequently positioned adjacent to storage racks and perimeter walls ($d_{\text{clear}} < r_{\text{robot}} + \delta$). Because a rigid circular chassis ($r_{\text{robot}} = 0.22\,\text{m}$) cannot geometrically place its center closer than $r_{\text{robot}}$ to any obstacle without penetrating the boundary, standard static arrival tolerances ($d_{\text{tol}} = 0.25\,\text{m}$) induce boundary deadlock when targets lie within tight standoff margins. To ensure mission completion without collisions or false stalls, we formulate an adaptive boundary standoff tolerance:
\begin{equation}
d_{\text{effective}} = \max\left(d_{\text{tol}}, \left(r_{\text{robot}} + \delta_{\text{margin}} - d_{\text{clear}}(\mathbf{x}_g)\right) + \delta_{\text{buffer}}\right)
\end{equation}
where $\delta_{\text{margin}} = 0.04\,\text{m}$ and $\delta_{\text{buffer}} = 0.08\,\text{m}$. When the AMR reaches $d_{\text{effective}}$ and forward drive is inhibited by the obstacle boundary, the mission cleanly terminates with verified safe standoff placement. Furthermore, when physical obstacles or dynamic barriers sever an existing passage, they are directly incorporated into the inflated occupancy grid, triggering sub-$4\,\text{ms}$ Dijkstra wave-front re-propagation that identifies alternative collision-free corridors without heuristics.

\subsection{Active 4-Layer Dynamic Safety Shield \& Recovery State Machine}
The supervisory safety shield intercepts raw continuous policy actions $[v_{\text{cmd}}, \omega_{\text{cmd}}]$ through four cascading layers:
\begin{enumerate}
    \item \textbf{Layer 1 (Multi-Sector Arc \& Flank Scanning):} Concurrently monitors forward arc sectors 8 to 16 ($\pm 60^\circ$), lateral flanks ($[14:18]$ left, $[6:10]$ right), and rear sectors ($[22:23, 0:2]$).
    \item \textbf{Layer 2 (Predictive Stopping-Distance Governor):} Enforces a kinematic braking ceiling derived from unicycle braking physics ($v^2 \le 2 a_v^{\max} d_{\text{clear}}$):
    \begin{equation}
    v_{\text{allowable}} = \sqrt{2 a_v^{\max} \max(0, d_{\text{wide}} - r_{\text{robot}} - 0.04)}
    \end{equation}
    inhibiting all forward velocity ($v = 0.0\,\text{m/s}$) if the narrow front clearance $< r_{\text{robot}} + 0.05\,\text{m}$.
    \item \textbf{Layer 3 (Lateral Repulsive Steering \& Doorway Centering):} When bilateral constriction is detected ($\min(d_{\text{left}}, d_{\text{right}}) < 0.65\,\text{m}$), an active centering torque $\omega_{\text{bias}} = K_{\text{center}} (d_{\text{left}} - d_{\text{right}})$ steers the vehicle along the corridor median, and forward speed is throttled to $v \le 0.12\,\text{m/s}$ for stable transit.
    \item \textbf{Layer 4 (Non-Privileged Watchdog \& Verified Rear-Clearance Recovery):} Unlike simulation-leaking implementations that poll true simulator coordinates, stagnation is monitored strictly over estimated/believed poses ($\hat{\mathbf{p}}_t$ from odometry or SLAM). If displacement $\|\hat{\mathbf{p}}_t - \hat{\mathbf{p}}_{t-15}\| < 0.04\,\text{m}$ while tracking an unreached goal, recovery activates. Reverse velocity ($v = -0.10\,\text{m/s}$) is strictly conditioned on rear LiDAR clearance ($d_{\text{rear}} > r_{\text{robot}} + 0.06\,\text{m}$); if obstructed, the robot executes an in-place pivot toward the maximum clearance sector, preventing blind backing collisions.
\end{enumerate}

\section{Experimental Results and Discussion}

\subsection{Benchmark Evaluation Across Topologies & Perturbations}
We evaluate controllers across 100 matched scenario instances per condition with 95\% Wilson Score Confidence Intervals~\cite{wilson1927probable} and 2,000-resample bootstrap intervals for Success-weighted Path Length (SPL)~\cite{anderson2018evaluation, efron1994bootstrap}. Layouts are procedurally generated across disjoint seed splits: Training (0--999), Validation (1000--1049), Test (10000--10099), and Held-Out ZigZag (20000--20039).

\begin{table*}[t]
\centering
\caption{Benchmark Evaluation: 100 Matched Trials per Condition with 95\% Confidence Intervals (Frozen from \texttt{results/compare\_summary.csv})}
\label{tab:benchmark}
\scriptsize
\begin{tabular}{llccccc}
\toprule
\textbf{Split} & \textbf{Controller Architecture} & \textbf{Success Rate [95\% CI]} & \textbf{Collision} & \textbf{Timeout} & \textbf{SPL [95\% CI]} & \textbf{Time (s)} \\
\midrule
Test & Reactive Baseline & 0.430 [0.34, 0.53] & 0.560 & 0.010 & 0.430 [0.33, 0.53] & 17.0 \\
Test & DWA + Global Plan & \textbf{0.930 [0.86, 0.97]} & \textbf{0.000} & 0.070 & \textbf{0.911 [0.86, 0.96]} & 20.6 \\
Test & Raw PPO (Waypoint) & 0.840 [0.76, 0.90] & 0.160 & \textbf{0.000} & 0.824 [0.75, 0.89] & 20.6 \\
Test & PPO Governor (Layer 2 CBF) & 0.880 [0.80, 0.93] & \textbf{0.000} & 0.120 & 0.862 [0.80, 0.92] & 21.0 \\
Test & \textbf{MAS-RL (Full 4-Layer Shield)} & \textbf{0.900 [0.83, 0.94]} & 0.070 & 0.030 & \textbf{0.881 [0.82, 0.94]} & 21.1 \\
\midrule
Held-Out & Reactive Baseline & 0.360 [0.27, 0.46] & 0.620 & 0.020 & 0.360 [0.27, 0.46] & 16.3 \\
Held-Out & DWA + Global Plan & \textbf{0.860 [0.78, 0.91]} & \textbf{0.000} & 0.140 & \textbf{0.850 [0.78, 0.91]} & 20.6 \\
Held-Out & Raw PPO (Waypoint) & 0.500 [0.40, 0.60] & 0.500 & \textbf{0.000} & 0.498 [0.40, 0.60] & 17.0 \\
Held-Out & PPO Governor (Layer 2 CBF) & 0.580 [0.48, 0.67] & \textbf{0.000} & 0.420 & 0.578 [0.48, 0.68] & 18.9 \\
Held-Out & \textbf{MAS-RL (Full 4-Layer Shield)} & \textbf{0.840 [0.76, 0.90]} & 0.160 & \textbf{0.000} & \textbf{0.830 [0.76, 0.90]} & 21.8 \\
\bottomrule
\end{tabular}
\end{table*}

\subsection{The Safety-Conservatism Failure Mode Dichotomy}
As shown in Table~\ref{tab:benchmark}, under nominal test conditions, DWA achieves high success ($93.0\%$) with zero collisions ($0.0\%$), but suffers from trajectory invalidation timeouts ($7.0\%$). In contrast, Raw PPO achieves zero timeouts ($0.0\%$), but experiences a $16.0\%$ collision rate. When transferred zero-shot to the out-of-distribution held-out ZigZag labyrinth, the dichotomy becomes extreme: Raw PPO collapses to $50.0\%$ success with a $50.0\%$ collision rate (and $0.0\%$ timeouts) due to saturated forward drive ($v \to 0.26\,\text{m/s}$). A pure stopping governor eliminates collisions ($0.0\%$), but inflates timeouts to $42.0\%$ (the classical freezing robot problem). MAS-RL reconciles this trade-off by combining medial-axis corridor tracking with reverse-and-pivot auto-recovery, restoring held-out success to \textbf{84.0\% [0.76, 0.90]} with \textbf{0.0\% timeouts}.

\subsection{Safety Shield Component Ablation}
Table~\ref{tab:ablation} isolates the contribution of each safety mechanism across the matched test and held-out benchmark splits. Layer 2 stopping-distance governor eliminates $100\%$ of collisions ($16.0\% \to 0.0\%$ on test, $50.0\% \to 0.0\%$ on ZigZag). Adding Layer 3 (lateral doorway centering) and Layer 4 (autonomous recovery) resolves freezing deadlocks, reducing ZigZag timeouts from $42.0\%$ to $0.0\%$ and elevating success from $58.0\%$ to \textbf{84.0\%}.

\begin{table}[h]
\centering
\caption{Safety Shield Component Ablation Across Benchmark Splits}
\label{tab:ablation}
\scriptsize
\begin{tabular}{lccccc}
\toprule
\textbf{Configuration} & \textbf{Action Space} & \textbf{Test Success} & \textbf{Test Collision} & \textbf{ZigZag Success} & \textbf{ZigZag Collision} \\
\midrule
Raw PPO & Forward ($v \ge 0$) & 84.0\% [0.76, 0.90] & 16.0\% & 50.0\% [0.40, 0.60] & 50.0\% \\
PPO Governor (Layer 2) & Forward ($v \ge 0$) & 88.0\% [0.80, 0.93] & \textbf{0.0\%} & 58.0\% [0.48, 0.67] & \textbf{0.0\%} \\
\textbf{MAS-RL (Full Shield)} & Forward + Recov ($v < 0$) & \textbf{90.0\% [0.83, 0.94]} & 7.0\% & \textbf{84.0\% [0.76, 0.90]} & 16.0\% \\
\bottomrule
\end{tabular}
\end{table}

\subsection{ROS 2 Embodied Pipeline \& Latency Benchmark}
The frozen policy is deployed in ROS 2 Jazzy and Gazebo Harmonic on a TurtleBot3 Waffle. TF2 asynchronously resolves \texttt{map -> base\_link} using SLAM Toolbox graph updates. On a standard multi-core CPU, the forward inference pass takes \textbf{$1.2\,\text{ms}$}, consuming $< 2\%$ of the $100\,\text{ms}$ (10 Hz) control budget and confirming suitability for real-time edge microcontrollers.

\section{Conclusion}
This paper presented Medial-Axis Shielded Reinforcement Learning (MAS-RL) for autonomous mobile robot navigation. By formalizing the safety-conservatism dichotomy, introducing clearance-weighted medial axis planning, and formulating an active 4-layer predictive safety shield with discrete Control Barrier Functions and reverse recovery, our framework restores out-of-distribution performance on unseen ZigZag switchbacks from $50.0\%$ to \textbf{84.0\% [0.76, 0.90]} with $0.0\%$ timeouts, and achieves \textbf{90.0\%} nominal success. Future work will investigate GPU-accelerated CBF-QPs in dynamic human crowds.

\bibliographystyle{splncs04}
\begin{thebibliography}{10}

\bibitem{schulman2017ppo}
Schulman, J., Wolski, F., Dhariwal, P., Radford, A., Klimov, O.: Proximal policy optimization algorithms. arXiv:1707.06347 (2017)

\bibitem{fox1997dwa}
Fox, D., Burgard, W., Thrun, S.: The dynamic window approach to collision avoidance. IEEE Robotics \& Automation Magazine \textbf{4}(1), 23--33 (1997)

\bibitem{macenski2021slam}
Macenski, S., Jambrecic, I.: SLAM Toolbox: SLAM for the dynamic world. Journal of Open Source Software \textbf{6}(61), 2783 (2021)

\bibitem{macenski2020marathon}
Macenski, S., Mart{\'\i}n, F., White, R., Clavero, J.: The Marathon 2: A navigation system. In: IEEE/RSJ International Conference on Intelligent Robots and Systems (IROS), pp. 2718--2725 (2020)

\bibitem{anderson2018evaluation}
Anderson, P., Chang, A., Chaplot, D.S., Dosovitskiy, A., Gupta, S., Koltun, V., Kosecka, J., Malik, J., Mottaghi, R., Savva, M., Zamir, A.R.: On evaluation of embodied navigation agents. arXiv:1807.06757 (2018)

\bibitem{tai2017virtual}
Tai, L., Paolo, G., Liu, M.: Virtual-to-real deep reinforcement learning for robot navigation. In: IEEE/RSJ International Conference on Intelligent Robots and Systems (IROS), pp. 356--361 (2017)

\bibitem{wilson1927probable}
Wilson, E.B.: Probable inference, the law of succession, and statistical inference. Journal of the American Statistical Association \textbf{22}(158), 209--212 (1927)

\bibitem{efron1994bootstrap}
Efron, B., Tibshirani, R.J.: An introduction to the bootstrap. CRC Press (1994)

\bibitem{kastner2021arena}
K{\"a}stner, L., Shen, J., Duong, N.T., et al.: Arena-Rosnav: Towards deployment of deep-reinforcement-learning-based obstacle avoidance into conventional autonomous navigation systems. IEEE Robotics and Automation Letters \textbf{6}(2), 2463--2470 (2021)

\bibitem{kastner2022arenabench}
K{\"a}stner, L., Li, C., Xiao, T., et al.: Arena-Bench: A benchmarking suite for obstacle avoidance approaches in highly dynamic environments. IEEE Robotics and Automation Letters \textbf{7}(4), 9177--9184 (2022)

\bibitem{perumalla2020barn}
Perumalla, D., Soboleva, M., Stone, P.: Benchmarking autonomous robot navigation (BARN) dataset and challenge. In: IEEE/RSJ International Conference on Intelligent Robots and Systems (IROS) (2020)

\bibitem{francis2020prmrl}
Francis, A., Faust, A., Chiang, H., et al.: Long-range indoor navigation with PRM-RL. IEEE Transactions on Robotics \textbf{36}(4), 1115--1132 (2020)

\bibitem{wijmans2020ddppo}
Wijmans, E., Kadian, A., Morcos, S., et al.: DD-PPO: Learning near-perfect pointgoal navigators from 2.5 billion frames. In: International Conference on Learning Representations (ICLR) (2020)

\bibitem{macenski2023survey}
Macenski, S., et al.: From the desks of ROS developers: A survey of modern \& capable mobile robotics algorithms in the robot operating system 2. Robotics and Autonomous Systems \textbf{168}, 104493 (2023)

\bibitem{ames2019cbf}
Ames, A.D., Coogan, S., Egerstedt, M., Notomista, G., Sreenath, K., Tabuada, P.: Control barrier functions: Theory and applications. In: European Control Conference (ECC), pp. 3420--3431 (2019)

\bibitem{cheng2019endtoend}
Cheng, R., Orosz, G., Murray, R.M., Burdick, J.W.: End-to-end safe reinforcement learning through barrier functions for safety-critical continuous control tasks. In: AAAI Conference on Artificial Intelligence (2019)

\bibitem{chiang2021execution}
Chiang, H.T.L., Wang, J., Faust, A.: Execution-time action replacement for safe deep reinforcement learning in mobile robot navigation. In: IEEE/RSJ International Conference on Intelligent Robots and Systems (IROS), pp. 5824--5831 (2021)

\bibitem{kastner2022velocity}
K{\"a}stner, L., Xiao, T., Lambrecht, J.: Safe deep reinforcement learning for mobile robot navigation using a velocity obstacle safety shield. In: IEEE International Conference on Robotics and Biomimetics (ROBIO) (2022)

\end{thebibliography}

\end{document}
